\chapter{Grundlagen}
\label{grundlagen}
Im Folgenden werden die theoretischen Grundlagen der Authentifizierungs- und Autorisierungsprotokolle in Abschnitt \ref{sec:authen_grundlagen} erklärt, dabei wird jeweils die Umsetzung in \ac{AWS} Cognito und \ac{SAP IAS} angedeutet. Des Weiteren wird in Abschnitt \ref{sec:tech_frameworks} die Technologische Integration der \ac{IdP}s vorgestellt, deren Architektur und Design-Prinzipien gegenübergestellt.

\section{Theoretische Grundlagen der Authentifizierungs- und Autorisierungsprotokolle}
\label{sec:authen_grundlagen}

Die Authentifizierung und Autorisierung von Benutzern in verteilten Cloud-Systemen basiert auf etablierten Standards und Protokollen, die eine sichere und interoperable Identitätsverwaltung ermöglichen. Im Rahmen dieser Arbeit werden die drei zentralen Protokolle OAuth 2.0, OpenID Connect und SAML 2.0 betrachtet, da sie die Grundlage für moderne \ac{IdP} wie \ac{AWS} Cognito und \ac{SAP IAS} bilden. Diese Protokolle unterscheiden sich in ihren Anwendungsfällen, ihrer technischen Umsetzung und den Sicherheitsmechanismen.

\subsection{OAuth 2.0 als Autorisierungsframework für delegierten Zugriff}

OAuth 2.0 ist ein Autorisierungsframework, da es Anwendungen ermöglicht, im Namen eines Benutzers auf geschützte Ressourcen zuzugreifen, ohne dass die Anmeldedaten des Benutzers mit der Anwendung geteilt werden müssen\footnote{ \url{https://datatracker.ietf.org/doc/html/rfc6749}, aufgerufen am 07.11.2025}. Das Framework definiert dabei vier Hauptakteure: den Resource Owner (z.B. der Benutzer), den Client (z.B. eine Web App), den Authorization Server (der die Authentifizierung durchführt und Token ausstellt) und den Resource Server (der die geschützten Ressourcen besitzt). 
\begin{figure}[h]
  \centering
  \includegraphics[width=1.05\textwidth]{figures/oauth2flow.drawio.pdf}
  \caption{Vereinfachter OAuth 2.0 Authorization Code Grant Flow}
  \label{fig:seq}
\end{figure}
Das Sequenzdiagramm \ref{fig:seq} stellt einen beispielhaften Authorization Code Grant Flow mit OAuth 2.0 dar.
Dabei beginnt es bei dem Resource Owner, dieser führt eine Aktion aus, die die Autorisierung startet, wie z.B. ein Login. Der Client leitet zunächst die Anfrage zum Authorization Server weiter. Der Authorization Server liefert dabei einen Endpunkt, wo sich der Resource Owner autorisieren kann, z.B. ein Login-Formular. Danach authentifiziert sich der Resource Owner gegenüber dem Authorization Server und der Authorization Server vergibt z.B. eine Genehmigung. Anschließend sendet der Authorization Server einen Authorization Code an den Client. Der Client tauscht diesen Code direkt beim Authorization Server gegen ein Access Token aus. Der Client nutzt 
schließlich den Access Token, um auf geschützte Ressourcen beim Resource-Server
 zuzugreifen. Der Vorteil liegt darin, dass der Access Token 
nicht über den Browser des Benutzers übertragen wird, sondern nur 
über den Client und dem Resource Server.
Der grundlegende OAuth 2.0-Ablauf folgt einem Request-Response-Muster mit unterschiedlichen \textit{Grant Types} für verschiedene Szenarien. Grant Types werden von OAuth 2.0 definiert für verschiedene Anwendungsszenarien.
Der \textit{Authorization Code Grant} ist der empfohlene Flow für die meisten Anwendungstypen und bietet die höchste Sicherheit, der der Access Token niemals über den Browser des Benutzers übertragen wird, sondern nur über Server-zu-Server-Kommunikation ausgetauscht wird. Beispielsweise wird bei einem Angular Frontend mit einem Java Backend als Server der Token mit dem Java Backend ausgetauscht.
Der \textit{Implicit Grant} wurde ursprünglich für browserbasierte Anwendungen, wie z.B. \ac{SPA} entwickelt, ist jedoch heute als unsicher klassifiziert und sollte nicht mehr verwendet werden. Für moderne \ac{SPA} wird stattdessen der Authorization Code Flow mit \ac{PKCE} empfohlen, der die Sicherheit beibehält und für Browser-Clients, wie unter anderem Angular, geeignet ist.
Der \textit{Resource Owner Password Credentials Grant} ist ein Flow, bei dem die Anwendung direkt die Anmeldedaten des Resource Owners erhält. Dieser Flow sollte daher nur in vertrauenswürdigen Umgebungen verwendet werden.
Der \textit{Client Credentials Grant} wird für Machine-to-Machine-Kommunikation ohne Benutzerinteraktion, bei der sich der Client direkt beim Authorization Server mit seinen eigenen Credentials authentifiziert verwendet.

OAuth 2.0 definiert zudem das Konzept der Scopes, die den Umfang der erteilten Berechtigung spezifizieren. Ein Scope beschränkt die Zugriffsrechte des Clients auf bestimmte Ressourcen oder Operationen. Der Authorization Server kann die vom Client angeforderten Scopes vollständig oder teilweise gewähren, basierend auf der Konfiguration. Access Tokens sind zeitlich begrenzt und enthalten Informationen über die gewährten Berechtigungen. Zusätzlich können Refresh Tokens ausgestellt werden, die dem Client ermöglichen, neue Access Tokens zu erhalten, ohne dass eine erneute Benutzerinteraktion erforderlich ist.

\subsubsection{Umsetzung in AWS Cognito}

AWS Cognito User Pools implementieren OAuth 2.0 als zentrales Autorisierungsframework und unterstützen alle gängigen Grant Types\footnote{\url{https://aws.amazon.com/blogs/security/how-to-use-oauth-2-0-in-amazon-cognito-learn-about-the-different-oauth-2-0-grants/}, aufgerufen am 07.11.2025}. Die Implementierung erfolgt über User Pools, die als Authorization Server fungieren und OAuth 2.0-konforme Tokens ausstellen. Cognito unterstützt den Authorization Code Flow mit \ac{PKCE} für erhöhte Sicherheit bei öffentlichen Clients, den Client Credentials Flow für Service-zu-Service-Kommunikation sowie die Verknüpfung mit externen \ac{IdP}s über OAuth 2.0. 

Die Konfiguration von OAuth 2.0 in Cognito erfolgt über Application Clients, die pro Mandant oder Anwendung erstellt werden können. Jede Client-Anwendung erhält eine eindeutige ID und optional ein Client Secret. Cognito ermöglicht die Konfiguration von OAuth-Scopes, die den Zugriff auf Benutzerattribute und API-Endpunkte steuern. Die Tokens werden im \ac{JWT}-Format ausgestellt und enthalten Claims über die Benutzeridentität und die gewährten Berechtigungen.

\subsubsection{Umsetzung in SAP IAS mit SAP CAP}

\ac{SAP IAS} implementiert OAuth 2.0 als Teil seiner Authentifizierungs- und Autorisierungsinfrastruktur und integriert sich nahtlos mit dem \ac{SAP CAP}-Framework\footnote{\url{https://cap.cloud.sap/docs/guides/security/}, aufgerufen am 07.11.2025}. Die Autorisierung wird dabei durch \ac{SAP CAP} verwaltet, während \ac{SAP IAS} die Authentifizierung übernimmt. \ac{SAP CAP} nutzt OAuth 2.0 Access Tokens zur Autorisierung von API-Anfragen und implementiert eine rollenbasierte Zugriffskontrolle, die auf den im Token enthaltenen Scopes basiert.

Die Integration erfolgt über die Konfiguration von OAuth 2.0-Clients in \ac{SAP IAS}, die mit \ac{SAP CAP}-Anwendungen verbunden werden. \ac{SAP CAP} bietet deklarative Sicherheitsannotationen, mit denen Entwickler Zugriffsbeschränkungen direkt im Datenmodell definieren können. Die Validierung der Access Tokens und die Durchsetzung der Autorisierungsregeln erfolgen automatisch durch das Framework. \ac{SAP IAS} unterstützt zudem die Verknüpfungen mit externen \ac{IdP}s über OAuth 2.0 und OpenID Connect.

\subsection{OpenID Connect als Authentifizierungsschicht auf OAuth 2.0-Basis}

OpenID Connect ist eine Authentifizierungsschicht, die auf OAuth 2.0 aufbaut und die standardisierte Überprüfung der Benutzeridentität ermöglicht\footnote{\url{https://openid.net/specs/openid-connect-core-1_0.html}, aufgerufen am 07.11.2025}. Während OAuth 2.0 die Autorisierung behandelt, erweitert \ac{OIDC} das Framework um Authentifizierungsfunktionen, indem es einen neuen Token-Typ einführt, nämlich einen ID-Token. Dieser Token enthält Claims über die Authentifizierung des Benutzers und ermöglicht es Client-Anwendungen, die Identität des Benutzers zu verifizieren und grundlegende Profilinformationen abzurufen.

Der \ac{OIDC}-Protokollablauf ähnelt den OAuth 2.0 Flows, verwendet jedoch den zusätzlichen Scope \texttt{openid}, um die Ausstellung eines ID Tokens zu signalisieren. Der ID-Token ist, wie der Access Token, ein \ac{JWT}, das signiert und optional verschlüsselt wird. Es enthält standardisierte Claims, wie z.B. \texttt{iss} (Issuer Identifier), \texttt{sub} (Benutzerkennung), \texttt{aud} (Audience oder Client ID), \texttt{exp} (Expiration Time), \texttt{iat} (Issued At Time) und \texttt{auth\_time} (Zeitpunkt der Authentifizierung). Zusätzlich können spezifische Claims wie E-Mail-Adresse, Name oder Rollen hinzugefügt werden.

\ac{OIDC} definiert drei primäre Authentifizierungs-Flows,zu einem dem \textit{Authorization Code Flow}, der auch der empfohlene Flow für die meisten Anwendungstypen ist, da er die höchste Sicherheit bietet. Der Authorization Code wird über den Browser übertragen, während die Tokens über einen sicheren Back-Channel zwischen Client und Authorization Server ausgetauscht werden. Der \textit{Implicit Flow} wurde ursprünglich für \ac{SPA} entwickelt, gilt jedoch heute als unsicher und sollte durch den Authorization Code Flow mit \ac{PKCE} ersetzt werden. Der \textit{Hybrid Flow} kombiniert Elemente beider Flows.

Zudem ist ein Konzept von \ac{OIDC} ein Endpunkt für Benutzerinformationen, indem es Access Token als Bearer Token darstellt. Die Antwort enthält Claims über den Benutzer im JSON-Format. \ac{OIDC} unterstützt zudem Discovery-Mechanismen, bei denen der OpenID Provider seine Metadaten über einen standardisierten Well-Known-Endpunkt publiziert. Diese Metadaten enthalten Informationen über unterstützte Flows, Endpunkt-URLs, unterstützte Scopes und verwendete Signaturalgorithmen.

\subsubsection{Umsetzung in AWS Cognito}

AWS Cognito User Pools fungieren als vollständiger\ac{OIDC} und unterstützen alle standardisierten \ac{OIDC}-Flows\footnote{\url{https://docs.aws.amazon.com/cognito/latest/developerguide/cognito-user-pools-oidc-flow.html}, aufgerufen am 07.11.2025}. Die Implementierung umfasst die automatische Ausstellung von ID-Tokens nach erfolgreicher Authentifizierung, einen Benuterinformations-Endpunkt zur Abfrage von Benutzerattributen sowie einen Discovery-Endpunkt für die automatische Konfiguration von Clients. Cognito stellt dabei drei Token-Typen aus. Zu einem dem ID-Tokens zur Authentifizierung, Access Tokens zur Autorisierung und Refresh Tokens zur Verlängerung von Sessions.

Die ID-Tokens werden als signierte \ac{JWT} ausgestellt und enthalten sowohl standardisierte als auch benutzerdefinierte Claims. Cognito verwendet RSA-Signaturalgorithmen (RS256) zur Signierung der Tokens und publiziert die öffentlichen Schlüssel über einen \ac{JWKS}-Endpunkt. Client-Anwendungen können die Token-Signatur validieren, um die Authentizität und Integrität der Tokens zu überprüfen. Cognito unterstützt zudem die Verknüpfung mit externen \ac{OIDC}-kompatiblen \ac{IdP}, was die Integration von Social Logins (Google, Facebook) oder Unternehmens-IdPs ermöglicht.

Die Konfiguration erfolgt über die Hosted UI von Cognito, die eine vollständig verwaltete Anmeldeoberfläche bereitstellt, oder über benutzerdefinierte UI-Implementierungen, die direkt mit den Cognito-APIs interagieren. Cognito ermöglicht die Anpassung der ausgestellten Claims durch serverlose Methoden, die vor der Token-Generierung ausgeführt werden können.

\subsubsection{Umsetzung in SAP IAS mit SAP CAP}

\ac{SAP IAS} ist als OpenID Connect-kompatibler \ac{IdP} konzipiert und integriert sich mit \ac{SAP CAP}-Anwendungen\footnote{\url{https://help.sap.com/docs/btp/sap-business-technology-platform/migration-from-saml-trust-to-openid-connect-trust-with-identity-authentication}, aufgerufen am 07.11.2025}. Die Implementierung folgt den \ac{OIDC} Core 1.0-Spezifikationen und unterstützt den Authorization Code Flow als primären Authentifizierungsmechanismus. \ac{SAP IAS} stellt ID-Tokens, Access Tokens und Refresh Tokens gemäß den \ac{OIDC}-Standards aus und ermöglicht die Abfrage von Benutzerattributen über den Bernutzerinformationen-Endpunkt.

Die Integration mit \ac{SAP CAP} erfolgt durch die Konfiguration von OAuth 2.0-Clients in \ac{SAP IAS} und die Registrierung der \ac{SAP CAP}-Anwendung. \ac{SAP CAP} nutzt die von \ac{SAP IAS} ausgestellten ID-Tokens zur Authentifizierung und die Access Tokens zur Autorisierung von API-Anfragen. Das Framework validiert automatisch die Token-Signatur und extrahiert die enthaltenen Claims für die Zugriffskontrolle. \ac{SAP IAS} unterstützt zudem die Anpassung der ausgestellten Claims über Attribute Mappings, die eine Transformation von Benutzerattributen während des Authentifizierungsvorgangs ermöglichen.

Die Konfiguration von \ac{SAP IAS} als \ac{OIDC}-Provider erfolgt über die Administration-Konsole, in der Anwendungen registriert, Redirect-URIs konfiguriert und Client-Credentials verwaltet werden. \ac{SAP IAS} bietet erweiterte Sicherheitsfeatures wie die Unterstützung von \ac{PKCE} für öffentliche Clients, die Konfiguration von Token-Gültigkeitsdauern und die Möglichkeit zur Anbindung externer \ac{IdP} über \ac{OIDC}-Verknüpfungen. Die Migration von bestehenden \ac{SAML}-basierten Integrationen zu \ac{OIDC} wird auch unterstützt, was die Modernisierung bestehender Landschaften erleichtert.

\subsection{SAML 2.0 - XML-basierter Standard für Enterprise Single Sign-On}

 \ac{SAML} 2.0 ist ein XML-basierter Standard für den Austausch von Authentifizierungs- und Autorisierungsinformationen zwischen vertrauenswürdigen Parteien\footnote{\url{https://docs.oasis-open.org/security/saml/Post2.0/sstc-saml-tech-overview-2.0.html}, aufgerufen am 07.11.2025}. \ac{SAML} wurde von der OASIS-Gruppe entwickelt und ist besonders im Enterprise-Umfeld für \ac{SSO}-Szenarien weit verbreitet. Das Protokoll definiert drei Hauptakteure. Der Principal (üblicherweise der Endbenutzer), den \ac{IdP} der die Authentifizierung durchführt und der Service Provider, der die Anwendung oder den Dienst bereitstellt, auf den zugegriffen werden soll.

Der zentrale Bestandteil von  \ac{SAML} ist die \textit{Assertion}, ein XML-Dokument, das Aussagen über einen Principal enthält. \ac{SAML} definiert drei Arten von Statements.
Die \textit{Authentication Statements} bestätigen, dass ein Principal zu einem bestimmten Zeitpunkt mit einer bestimmten Methode authentifiziert wurde. \textit{Attribute Statements} enthalten Attributinformationen über den Principal, wie Name, E-Mail-Adresse oder Gruppen. \textit{Authorization Decision Statements} beschreiben, ob ein Principal berechtigt ist, eine bestimmte Aktion auf einer Ressource auszuführen. Die Assertions werden signiert, um deren Integrität und Authentizität zu gewährleisten, und können optional verschlüsselt werden.

Der \ac{SAML}-Authentifizierungsablauf im Web Browser \ac{SSO}-Profil beginnt damit, dass der Benutzer versucht auf eine geschützte Ressource beim Service Provider zuzugreifen. Der Service Provider erkennt, dass der Benutzer nicht authentifiziert ist und leitet ihn zum konfigurierten \ac{IdP} weiter. Der \ac{IdP} authentifiziert den Benutzer (z.B. über Username/Passwort oder Multi-Faktor-Authentifizierung) und erstellt eine \ac{SAML} Assertion, die die Identität des Benutzers bestätigt. Die Assertion wird signiert und über den Browser des Benutzers an den Service Provider übermittelt. Der Service Provider validiert die Signatur der Assertion und gewährt basierend auf den enthaltenen Informationen Zugriff auf die angeforderte Ressource.

\ac{SAML} definiert verschiedene Bindings, die festlegen, wie \ac{SAML}-Protokollnachrichten über verschiedene Transportprotokolle übertragen werden. Das HTTP Redirect Binding nutzt HTTP-Redirects zur Übertragung von \ac{SAML}-Requests, während das HTTP POST Binding SAML-Responses als POST-Parameter überträgt. Das SOAP Binding wird für direkte Service-zu-Service-Kommunikation verwendet. \ac{SAML}-Profiles kombinieren Assertions, Protocols und Bindings zu konkreten Anwendungsfällen, wobei das Web Browser \ac{SSO} Profile das am häufigsten verwendete ist.

\subsubsection{Umsetzung in AWS Cognito}

AWS Cognito User Pools unterstützen \ac{SAML} 2.0 sowohl als Service Provider als auch für die Verknüpfung mit externen \ac{SAML}-basierten \ac{IdP}\footnote{\url{https://docs.aws.amazon.com/cognito/}, aufgerufen am 07.11.2025}. Als Service Provider kann Cognito \ac{SAML} Assertions von Unternehmens-\ac{IdP}s wie Active Directory Services oder anderen \ac{SAML} 2.0-konformen \ac{IdP} akzeptieren. Die Integration ermöglicht es Unternehmen, bestehende Identitätsinfrastrukturen zu nutzen und Cognito als Gateway für den Zugriff auf AWS-Ressourcen und Anwendungen zu verwenden.

Die Konfiguration erfolgt durch den Import der \ac{SAML}-Metadaten des \ac{IdP} in Cognito und die Registrierung von Cognito als Service Provider beim \ac{IdP}. Cognito generiert automatisch \ac{SAML}-Metadaten, die an den \ac{IdP} übertragen werden müssen. Die Attribute Mapping-Funktion ermöglicht die Zuordnung von \ac{SAML}-Attributen zu Cognito-Benutzerattributen, sodass Informationen wie Name, E-Mail-Adresse oder Gruppenzugehörigkeiten automatisch synchronisiert werden. Nach erfolgreicher \ac{SAML}-Authentifizierung stellt Cognito OAuth 2.0/\ac{OIDC}-konforme Tokens aus, die für die Autorisierung von API-Anfragen verwendet werden können.

Cognito unterstützt zudem \ac{JIT} User Provisioning, bei dem Benutzer automatisch im User Pool angelegt werden, wenn sie sich zum ersten Mal über einen \ac{SAML}-\ac{IdP} anmelden. Die Attribute werden dabei aus der \ac{SAML} Assertion extrahiert und im Cognito-Benutzerprofil gespeichert.

\subsubsection{Umsetzung in SAP IAS mit SAP CAP}

\ac{SAP IAS} bietet auch Unterstützung für \ac{SAML} 2.0 und kann sowohl als \ac{IdP} als auch als Service Provider konfiguriert werden\footnote{\url{https://www.sap.com/products/technology-platform/identity-authentication.html}, aufgerufen am 07.11.2025}. Als \ac{IdP} kann \ac{SAP IAS} \ac{SAML} Assertions für SAP- und Nicht-SAP-Anwendungen ausstellen, die als Service Provider konfiguriert sind. Die Implementierung folgt den \ac{SAML} 2.0-Standards und unterstützt sowohl \ac{IdP}-initiated als auch Service Provider-initiated \ac{SSO}-Flows.

Die Integration mit \ac{SAP CAP}-Anwendungen erfolgt primär über \ac{OIDC}, jedoch kann \ac{SAP IAS} als Proxy zwischen \ac{SAML}-basierten Unternehmens-\ac{IdP} und \ac{OIDC}-basierten Anwendungen fungieren. Dies ermöglicht eine nahtlose Integration bestehender SAML-Infrastrukturen mit modernen \ac{SAP CAP}-Anwendungen. \ac{SAP IAS} unterstützt die Verknüpfung mit externen \ac{SAML}-\ac{IdP}s und kann als zentraler Hub für die Authentifizierung in SAP-Landschaften dienen.

Die Konfiguration umfasst den Import von \ac{SAML}-Metadaten externer \ac{IdP}s, die Definition von Attribute Mappings zur Transformation von \ac{SAML}-Attributen in \ac{SAP IAS}-Benutzerattribute und die Konfiguration von Trust Relationships zwischen \ac{SAP IAS} und Service Providern. \ac{SAP IAS} bietet erweiterte Features wie Conditional Authentication, die es ermöglicht, unterschiedliche Authentifizierungsrichtlinien basierend auf Kontextinformationen anzuwenden. Die Migration von \ac{SAML}-basierten zu \ac{OIDC}-basierten Integrationen wird durch SAP aktiv unterstützt, da \ac{OIDC} als zukunftssicherer Standard für moderne Cloud-Anwendungen betrachtet wird.

\subsection{Vergleich der Protokolle und Anwendungsszenarien}

Die drei Protokolle OAuth 2.0, OpenID Connect und \ac{SAML} 2.0 adressieren unterschiedliche Anwendungsfälle und weisen jeweils Stärken und Schwächen auf. OAuth 2.0 fokussiert sich auf die Autorisierung und delegierten Zugriff, ist jedoch kein Authentifizierungsprotokoll. Es eignet sich besonders für API-Zugriffe und Machine-to-Machine-Kommunikation. Die Verwendung von JSON-basierten Tokens und \ac{REST} APIs macht OAuth 2.0 entwicklerfreundlich und gut geeignet für z.B. Angular Web-Anwendungen.

OpenID Connect erweitert OAuth 2.0 um Authentifizierungsfunktionalität und ist das Protokoll der Wahl für moderne \ac{CIAM}-Anwendungen. Die Verwendung von \ac{JWT}-basierten ID Tokens, die Integration in Single-Page Applications und Mobile Apps sowie die Unterstützung für Social Logins machen \ac{OIDC} zur einer guten kundenorientierte Lösung für Anwendungen. Die JSON-basierte Kommunikation ist zudem leichter als XML, und die Protokollspezifikation ist auf moderne Cloud- und Mobile-Architekturen ausgerichtet.

\ac{SAML} 2.0 ist der etablierte Standard für Enterprise \ac{SSO} und wird primär in Unternehmensumgebungen eingesetzt. Die XML-basierte Struktur ermöglicht Attribute und Assertions, was \ac{SAML} besonders für Szenarien geeignet macht, in denen Benutzerinformationen und Autorisierungsentscheidungen übertragen werden. \ac{SAML} bietet Sicherheitsmechanismen durch XML-Signatur und -Verschlüsselung und ist gut in bestehende Enterprise-Infrastrukturen wie Active Directory integriert. Jedoch ist \ac{SAML} komplexer in der Implementierung und weniger geeignet für Single-Page-Anwendungen \citep{Zhuravel2024}.

Die Wahl des geeigneten Protokolls hängt von den spezifischen Anforderungen der Anwendung ab, denn für moderne z.B. Angular Web-Anwendungen mit Fokus auf Benutzerauthentifizierung ist \ac{OIDC} die bevorzugte Wahl. Für Enterprise-SSO-Szenarien mit Integration in bestehende \ac{SAML}-basierte Infrastrukturen bleibt \ac{SAML} 2.0 relevant. Für API-Zugriffe und Machine-to-Machine-Kommunikation ohne Benutzerinteraktion ist OAuth 2.0 ausreichend \citep{Zhuravel2024}.\ac{IdP}s wie AWS Cognito und \ac{SAP IAS} unterstützen alle drei Protokolle, um maximale Flexibilität und Interoperabilität zu gewährleisten.


\section{Technologische Integration der \ac{IdP} in die Anwendungs-Ökosysteme}
\label{sec:tech_frameworks}

Die praktische Integration der \ac{IdP}s erfolgt als Teil spezifischer Anwendungs-Frameworks. Das beeinflusst die Implementierungskomplexität, die Developer Experience und damit auch das Sicherheitsniveau. Dabei werden zwei unterschiedliche Szenarien miteinander verglichen, die beide Architekturen in der Praxis repräsentieren.

\subsection{\ac{SAP CAP}}

\ac{SAP CAP} ist ein Enterprise-Framework zur Entwicklung von Cloud-Anwendungen und bildet die Grundlage für die Integration mit \ac{SAP IAS}. Das Framework besteht aus verschiedenen Sprachen, Bibliotheken und Tools. Die \ac{SAP CAP} Service \ac{SDK} wird zudem für Node.js und Java angeboten, dabei wird der Fokus auf die Java Version gelegt. Zudem nutzt das \ac{SDK} \textbf{Spring Boot} als Grundlage für die technische Implementierung\footnote{\label{fn:sapcap}\url{https://cap.cloud.sap/docs/about/\#what-is-cap}, aufgerufen am 21.11.2025}.

\subsubsection{Architektur und Design-Prinzipien}

\ac{SAP CAP} folgt dem Prinzip der \textit{Domain-Driven Design} Entwicklung. Das Domain-Driven Design fokussiert sich darauf, jede Domäne in ihre eigenen Kontexte zu unterteilen und jeweils spezifische Modelle zugeordnet werden.\footnote{\url{https://cap.cloud.sap/docs/guides/domain-modeling}, aufgerufen 21.11.2025}. Das Framework verwendet \ac{CDS} als universelle Modellierungssprache zur Definierung von Datenmodellen, Services und Autorisierungslogik. Ein Design-Merkmal von \ac{SAP CAP} ist die strikte Trennung zwischen Datendefinition und Geschäftslogik, was eine deklarative Architektur ermöglicht \footnote{\ref{fn:sapcap}}.

Die Sicherheit in \ac{SAP CAP} wird primär durch Annotationen im \ac{CDS}-Modell konfiguriert. Dabei werden Autorisierungsregeln direkt im Datenmodell definiert mithilfe von Annotationen wie beispielsweise \texttt{@requires} und \texttt{@restrict}. Diese Regeln werden zur Laufzeit vom \ac{SAP CAP}-Framework automatisch durchgesetzt, ohne dass manuelle Validierungslogik im Service-Code erforderlich ist \footnote{\ref{fn:sapcap}}.

Das Framework integriert sich mit Spring Boot, aber es werden zusätzlich eigene Erweiterungen für die Authentifizierung und Autorisierung benutzt. Die automatische Spring Boot-Sicherheitskonfiguration wird durch die Bibliothek \texttt{cds-framework-spring-boot} bereitgestellt und aktiviert sich automatisch, wenn Abhängigkeiten und Identity-Service-Bindungen vorhanden sind.

\subsubsection{Integration mit \ac{SAP IAS}}

Die Integration von \ac{SAP IAS} in \ac{SAP CAP} erfolgt auf zwei Ebenen: Authentifizierung und Autorisierung. \ac{SAP IAS} übernimmt ausschließlich die Authentifizierung von Benutzern über standardisierte Protokolle (\ac{OIDC}/\ac{SAML}), während die Autorisierung vollständig von \ac{SAP CAP} verwaltet wird.

Der Autorisierungsmechanismus basiert auf deklarativen Annotationen. Beispielsweise können Entities mit einer Autorisierungsregel versehen:

\begin{verbatim}
entity Orders : cuid {
  customer: String;
  amount: Decimal;
}
annotate Orders with @(restrict: [
  { grant: 'READ,WRITE', to: 'Managers' }
]);
\end{verbatim}

Das \ac{SAP CAP}-Framework wertet diese Deklaration automatisch aus und erzeugt entsprechende SQL-Filter oder API-Zugriffsbeschränkungen. Dies bedeutet, dass nicht selbst Code geschrieben werden muss, um Autorisierungslogik zu implementieren. Dadurch wird auch das Risiko für Fehlkonfigurationen reduziert.

\subsection{Spring Boot mit AWS Cognito und Spring Security}

Für die Integration von \ac{AWS} Cognito in Java-Anwendungen ist Spring Boot mit Spring Security der etablierte Standard. Im Gegensatz zu \ac{SAP CAP} folgt dieser Ansatz einem imperativen Programmiermodell, bei dem Code für die Authentifizierung und Autorisierung geschrieben werden muss.

\subsubsection{Architektur und Design-Prinzipien}

Spring Boot ist ein Framework zur schnellen Entwicklung von eigenständigen Java-Anwendungen. Es folgt dem Prinzip \textit{Convention over Configuration} und vereinfacht die Ausführung und Verwaltung von Produktionsanwendungen. Spring Security ist eine spezialisierte Bibliothek innerhalb des Spring-Ökosystems, die flexible und konfigurierbare Sicherheitsmechanismen zur Verfügung stellt\footnote{\url{https://spring.io/}, aufgerufen am 21.11.2025}.

Der Sicherheitsansatz in Spring Boot/Spring Security basiert dabei auf Annotationen und Konventionen. Dabei werden Annotationen wie \texttt{@PreAuthorize}, \texttt{@Secured} oder \texttt{@RolesAllowed} auf Controller- oder Service-Methoden benutzt, um Zugriffsbeschränkungen zu definieren:

\begin{verbatim}
@RestController
@RequestMapping("/api/orders")
public class OrderController {
  @PreAuthorize("hasRole('MANAGER')")
  @GetMapping("/{id}")
  public Order getOrder(@PathVariable String id) { }
}
\end{verbatim}

Die Durchsetzung dieser Autorisierungsregeln erfolgt durch das Spring Security-Framework zur Laufzeit. Dies erfordert jedoch, dass diese Annotationen auf jeder Methode oder Klasse wiederholt werden müssen.

\subsubsection{Integration mit AWS Cognito}

\ac{AWS} Cognito wird durch Spring Security integriert, indem das Framework als \ac{OIDC}-Provider konfiguriert wird. Die Konfiguration erfolgt in der \texttt{application.yaml} oder \texttt{application.properties} Datei:

\begin{verbatim}
spring.security.oauth2.resourceserver.jwt.issuer-uri=
  https://cognito-idp.<region>.amazonaws.com/<user-pool-id>
spring.security.oauth2.resourceserver.jwt.jwk-set-uri=
  https://cognito-idp.<region>.amazonaws.com/<user-pool-id>/.well-known/jwks.json
\end{verbatim}

Nach dieser Konfiguration validiert Spring Security automatisch die \ac{JWT}-Tokens, die von \ac{AWS} Cognito ausgestellt werden. Die Autorisierungsregeln müssen jedoch manuell im Code implementiert werden. Zudem muss für jede geschützte Methode eine entsprechende Prüfung durchgeführt werden oder Annotationen verwenden werden.

Ein zentraler Unterschied zu \ac{SAP CAP} besteht darin, dass Spring Security keine deklarative Datenzugriffsschutzschicht bietet. Das bedeutet, dass für jeden Service manuell implementiert werden muss, welche Daten ein Benutzer mit bestimmten Rollen abrufen darf.

\subsubsection{Vergleich der Architektur-Ansätze}

Die beiden Integrationspfade unterscheiden sich in ihrem Design:

\begin{itemize}
    \item \textbf{Deklarativ vs. Imperativ:} \ac{SAP CAP} nutzt einen deklarativen Ansatz, bei dem Regeln im \ac{CDS}-Modell definiert werden. Spring Security/Spring Boot nutzt einen imperativen Ansatz, bei dem Entwickler Regeln an beliebigen Stellen im Code durch Annotationen oder Programmlogik implementieren.
    
    \item \textbf{Zentralisiert vs. Verteilt:} In \ac{SAP CAP} sind Autorisierungsregeln zentral im Datenmodell definiert. Bei Spring Security sind diese über Controller, Services und möglicherweise weitere Komponenten verteilt.
    
    \item \textbf{Framework-gesteuert vs. Developer-gesteuert:} \ac{SAP CAP} erzwingt eine bestimmte Struktur und Best Practice für Autorisierung. Spring Security bietet dafür Flexibilität.
    
    \item \textbf{Automatische Durchsetzung:} \ac{SAP CAP} führt Autorisierungsregeln automatisch auf Datenbankebene durch. Spring Security erfordert, dass jede Methode explizit geschützt wird.
\end{itemize}

Diese Unterschiede haben direkte Auswirkungen auf die Messbarkeit der in dieser Arbeit untersuchten Kriterien wie Sicherheit, Wartbarkeit und Developer Experience.


\subsection{Authentifizierungsschicht: Token-Austausch und Benutzerregistrierung}

Die Authentifizierungsschicht definiert, wie Benutzer sich anmelden und registrieren sowie wie das System Access Tokens ausstellt. Obwohl beide \ac{IdP} auf standardisierten Protokollen basieren, unterscheiden sich ihre Implementierungen.

\subsubsection{Passwort-Registrierung und Authentifizierung bei AWS Cognito}

Bei AWS Cognito muss sich jeder Benutzer mit Passwort registrieren. Das Passwort wird durch konfigurierbare Regeln (Mindestlänge, Groß-/Kleinbuchstaben, Zahlen, Sonderzeichen) erzwungen\footnote{\url{https://docs.aws.amazon.com/cognito/latest/developerguide/signing-up-users-in-your-app.html}, aufgerufen am 22.11.2025}. Nach Registrierung erfolgt eine E-Mail-Bestätigung.

Der Authentifizierungs-Flow folgt dem OAuth 2.0 Authorization Code Grant Flow, wie bereits in Abbildung \ref{fig:seq} beschrieben. Der Benutzer gibt seine Anmeldedaten ein, Cognito validiert diese und sendet einen Authorization Code zurück. Die Spring Boot-Anwendung tauscht diesen Code gegen \ac{JWT}-Tokens (ID Token, Access Token, Refresh Token) aus.
Cognito unterstützt mehrere Authentifizierungs-Flows. Zu einem wird der \ac{SRP}-Flow, der das Passwort nicht sendet sondern berechnet einen öffentliche Wert und sendet diesen, oder auch WebAuthn/Passkeys (biometrische oder Hardware-Key-Authentifizierung), One-Time-Passcodes via SMS oder Email, sowie adaptive Authentifizierung, die basierend auf Risikoscores dynamisch zusätzliche Authentifizierungsfaktoren erfordern kann, implementiert\footnote{\url{https://docs.aws.amazon.com/cognito/latest/developerguide/amazon-cognito-user-pools-authentication-flow-methods.html}, aufgerufen am 23.11.2025}.

Spring Security validiert diese Tokens automatisch über die Cognito-\ac{JWKS}-URI, wenn die Konfiguration vorgenommen wurde.

\subsubsection{Passwort-Registrierung und Authentifizierung bei \ac{SAP IAS}}

\ac{SAP IAS} besitzt ein ähnliches Konzept, dabei implementiert es standardisierte Passwort-Policies \textit{Tenant}-weit. Ein Tenant ist im SAP Kontext, eine isolierte zugewiesene Umgebung innerhalb einer gemeinsamen Cloud-Infrastruktur. Dadurch können beispielsweise Kunden zu unterschiedlichen Tenants zugewiesen werden\footnote{\url{https://help.sap.com/docs/SAP_Leonardo_IoT/080fabc6cae6423fb45fca7752adb61e/462b49382316427aa59fe671a75fa39e.html}, aufgerufen am 22.11.2025} Neue Benutzer können von Administratoren erstellt oder durch Self-Service registriert werden. Der Authentifizierungs-Flow ist ebenfalls \ac{OIDC}-basiert und folgt dem Authorization Code Flow (Abbildung \ref{fig:seq})\footnote{\url{https://help.sap.com/docs/cloud-identity-services/cloud-identity-services/create-new-user}, aufgerufen am 22.11.2025}.

Im Gegensatz zu Cognito unterstützt \ac{SAP IAS} auch den Resource Owner Password Credentials Flow, bei dem der Client Benutzerdaten direkt vom Benutzer erhält.
Außerdem bietet \ac{SAP IAS} als Alternativen zu Passwörtern X.509-Zertifikate sowie WebAuthn/Passkeys und Biometrie, unterstützt zudem Risk-Based Access Control (erhöhte Authentifizierungsanforderungen basierend auf Kontext wie Netzwerkstandort) und ermöglicht auch eine Integration mit externen \ac{IdP}\footnote{\url{https://www.erpqna.com/passwordless-authentication-passkeys-with-sap-btp-sdk-for-ios-and-sap-cloud-identity-services/}, aufgerufen am 23.11.2025}\footnote{\url{https://www.linkedin.com/pulse/strengthening-sap-authentication-zero-trust-methods-carsten-olt-dwiqe}, aufgerufen am 23.11.2025}.

In \ac{SAP CAP} erfolgt die Token-Validierung automatisch durch die \texttt{cds-framework-spring-boot}-Bibliothek. Das Framework konfiguriert Spring Security automatisch mit dem \ac{SAP IAS}-JWKS-Endpunkt, wenn eine Service-Binding vorhanden ist.

\paragraph{Unterschiede in der Authentifizierungsschicht}

\begin{itemize}
    \item \textbf{Token-Flows:} Cognito unterstützt primär Authorization Code Flow. \ac{SAP IAS} bietet mehr Flexibility mit mehreren Standards-Flows.
    
    \item \textbf{Passwort-Policies:} Cognito erlaubt mehrere Konfiguration (höhere Komplexität, Fehlerkonfigurationsrisiko). \ac{SAP IAS} standardisiert auf Best Practices (weniger Fehler, weniger Flexibilität).
    
    \item \textbf{Token-Validierung:} Cognito + Spring Boot erfordert manuelle Konfiguration. \ac{SAP CAP} + \ac{SAP IAS} validiert automatisch durch das Framework.
    
    \item \textbf{Developer Experience:} \ac{SAP CAP} reduziert durch Automatisierung die Komplexität. Spring Boot mit Cognito erfordert mehr Konfiguration.
\end{itemize}
